%% =========================================================
%% Chapter: Quantum Dynamical Semigroups and GKSL Generators
%% Covers: TNLean/Channel/Semigroup/{Basic,Perturbation,Primitivity,Generator}.lean
%% Wolf reference: Chapter 7
%% =========================================================
%% The semigroup declarations live in specialized modules that are not
%% re-exported from TNLean.lean, so this chapter intentionally records
%% formalization status without direct \lean{} tags.

\chapter{Quantum Dynamical Semigroups}
\label{ch:semigroup}

This chapter develops the theory of one-parameter semigroups of quantum
channels and characterizes their generators via the
GKSL (Gorini--Kossakowski--Sudarshan--Lindblad) theorem,
following~\cite[Ch.~7]{Wolf2012Quantum}.

Section~\ref{sec:semigroup_defs} establishes definitions and
the basic exponential form
(Proposition~\ref{thm:continuousDynSemigroup_eq_exp_ch12}).
Section~\ref{sec:perturbation} gives the Duhamel formula and
perturbation bound.
Section~\ref{sec:gksl_generators} develops the generator theory
and proves the GKSL theorem.

%% ---------------------------------------------------------
\section{Dynamical semigroups and the exponential form}
\label{sec:semigroup_defs}
%% ---------------------------------------------------------

\begin{definition}[Dynamical semigroup]\label{def:dyn_semigroup_ch12}
    \leanok
    A family of linear maps
    $T : \R_{\geq 0} \to (\MN{D} \to_{\ell} \MN{D})$
    is a \emph{dynamical semigroup} if
    \begin{enumerate}
        \item (semigroup law) $T(t+s) = T(t) \circ T(s)$
              for all $t,s \ge 0$; and
        \item (initial condition) $T(0) = \id$.
    \end{enumerate}
    This is \cite[Eq.~(7.1)]{Wolf2012Quantum}.
\end{definition}

\begin{definition}[Continuous dynamical semigroup]\label{def:continuous_dyn_semigroup_ch12}
    \leanok
    \uses{def:dyn_semigroup_ch12}
    A dynamical semigroup $T$ is \emph{norm-continuous} if the map
    $t \mapsto T(t)$ is continuous in the operator norm topology.
    In finite dimension this coincides with strong continuity.
\end{definition}

\begin{definition}[Exponential semigroup]\label{def:exp_semigroup_ch12}
    \leanok
    Given a generator $L \in \mathrm{End}_{\C}(\MN{D})$, the
    \emph{exponential semigroup} is
    \[
        T_t := e^{tL} = \sum_{k=0}^{\infty} \frac{(tL)^k}{k!}.
    \]
\end{definition}

\begin{theorem}[Semigroup law for the exponential]\label{thm:expSemigroupCLM_add_ch12}
    \leanok
    \uses{def:exp_semigroup_ch12}
    For all $t,s \in \R$,
    $e^{(t+s)L} = e^{tL} \cdot e^{sL}$.
\end{theorem}

\begin{proof}
    \leanok
    Since $(t \cdot L)$ commutes with $(s \cdot L)$, this follows from
    the exponential addition formula for commuting elements in a
    Banach algebra.
\end{proof}

\begin{theorem}[Continuity of the exponential semigroup]\label{thm:expSemigroupCLM_continuous_ch12}
    \leanok
    \uses{def:exp_semigroup_ch12}
    The map $t \mapsto e^{tL}$ is continuous in the operator norm topology.
\end{theorem}

\begin{proof}
    \leanok
    The exponential is analytic (convergence radius $= \infty$), hence
    continuous.
\end{proof}

\begin{theorem}[Derivative of the exponential semigroup]\label{thm:hasDerivAt_expSemigroupCLM_ch12}
    \leanok
    \uses{def:exp_semigroup_ch12}
    For all $t \in \R$,
    \[
        \frac{d}{dt} e^{tL} = e^{tL} \cdot L.
    \]
    This is \cite[Eq.~(7.2)]{Wolf2012Quantum}.
\end{theorem}

\begin{proof}
    \leanok
    Chain rule applied to the composition
    $\R \xrightarrow{t \mapsto t \cdot L} \mathrm{End}(\MN{D})
     \xrightarrow{\exp} \mathrm{End}(\MN{D})$.
\end{proof}

\begin{theorem}[Generator uniqueness]\label{thm:generator_unique_ch12}
    \leanok
    \uses{def:exp_semigroup_ch12}
    If $e^{tL} = e^{tL'}$ for all $t \ge 0$, then $L = L'$.
\end{theorem}

\begin{proof}
    \leanok
    Both semigroups agree on $[0,\infty)$, so their derivatives within
    $[0,\infty)$ at $t=0$ agree.
    Since $[0,\infty)$ has unique differentials at~$0$, $L = L'$.
\end{proof}

\begin{theorem}[Prop.~7.1: Every continuous semigroup is exponential]
    \label{thm:continuousDynSemigroup_eq_exp_ch12}
    \leanok
    \uses{def:continuous_dyn_semigroup_ch12, def:exp_semigroup_ch12,
          thm:generator_unique_ch12}
    Every norm-continuous dynamical semigroup $T$ on $\MN{D}$ is
    of the form
    \[
        T_t = e^{tL}
    \]
    for a unique generator $L \in \mathrm{End}_{\C}(\MN{D})$.
    This is \cite[Prop.~7.1]{Wolf2012Quantum}.
\end{theorem}

\begin{proof}
    \leanok
    Since $T_0 = \Id$ and $T$ is continuous, the integral
    $M_\varepsilon = \int_0^\varepsilon T_s\,ds$ is invertible for small
    $\varepsilon > 0$.
    Using the semigroup property,
    $T_t = M_\varepsilon^{-1}(M_{t+\varepsilon} - M_t)$
    is differentiable.
    ODE uniqueness gives $T_t = e^{tL}$ with $L = T'(0)$.
\end{proof}

%% ---------------------------------------------------------
\section{Perturbation theory}
\label{sec:perturbation}
%% ---------------------------------------------------------

\begin{theorem}[Lem.~7.1: Duhamel formula]\label{thm:duhamel_formula_ch12}
    \leanok
    \uses{def:exp_semigroup_ch12, thm:hasDerivAt_expSemigroupCLM_ch12}
    Let $\Delta = L' - L$. For $t \ge 0$,
    \[
        e^{tL'} - e^{tL}
          = \int_0^t e^{(t-s)L} \Delta\,e^{sL'}\,ds.
    \]
    This is \cite[Lem.~7.1]{Wolf2012Quantum}.
\end{theorem}

\begin{proof}
    \leanok
    Define $f(s) = e^{(t-s)L} e^{sL'}$.
    Then $f'(s) = e^{(t-s)L} \Delta\,e^{sL'}$
    (Theorem~\ref{thm:hasDerivAt_expSemigroupCLM_ch12}),
    and $e^{tL'} - e^{tL} = f(t) - f(0) = \int_0^t f'(s)\,ds$.
\end{proof}

\begin{theorem}[Cor.~7.1: Perturbation bound]\label{thm:perturbation_bound_ch12}
    \leanok
    \uses{thm:duhamel_formula_ch12}
    For $t \ge 0$,
    \[
        \bigl\|e^{tL'} - e^{tL}\bigr\|
          \le t\,\|\Delta\|
             \cdot \sup_{s \in [0,t]}\|e^{sL}\|
             \cdot \sup_{s \in [0,t]}\|e^{sL'}\|,
    \]
    where $\Delta = L' - L$.
    This is \cite[Cor.~7.1]{Wolf2012Quantum}.
\end{theorem}

\begin{proof}
    \leanok
    Apply the triangle inequality and submultiplicativity of the
    operator norm to the Duhamel integral
    (Theorem~\ref{thm:duhamel_formula_ch12}).
\end{proof}

%% ---------------------------------------------------------
\section{GKSL/Lindblad generators}
\label{sec:gksl_generators}
%% ---------------------------------------------------------

This section characterizes generators of semigroups of completely
positive and trace-preserving maps.
The central result is the GKSL theorem
(Theorem~\ref{thm:gksl_iff_lindblad}),
which shows that such generators have the standard Lindblad form.

\subsection{Generator decomposition and conditional complete positivity}

\begin{definition}[Generator decomposition]\label{def:generator_decomp}
    \leanok
    A \emph{generator decomposition} consists of a completely positive map
    $\phi : \MN{d} \to \MN{d}$ and a matrix $\kappa \in \MN{d}$, defining
    the linear map
    \[
        L(\rho) = \phi(\rho) - \kappa\rho - \rho\kappa^\dagger.
    \]
    This is \cite[Eq.~(7.14)]{Wolf2012Quantum}.
\end{definition}

\begin{definition}[Conditional complete positivity]\label{def:is_ccp}
    \leanok
    \uses{def:generator_decomp}
    A linear map $L : \MN{d} \to \MN{d}$ is \emph{conditionally
    completely positive} (CCP) if it admits a generator decomposition,
    i.e.\ there exist a CP map $\phi$ and a matrix $\kappa$ such that
    $L(\rho) = \phi(\rho) - \kappa\rho - \rho\kappa^\dagger$.
    This is condition~1 of \cite[Prop.~7.2]{Wolf2012Quantum}.
\end{definition}

\begin{definition}[Trace-annihilating map]\label{def:is_trace_annihilating}
    \leanok
    A linear map $L : \MN{d} \to \MN{d}$ is \emph{trace-annihilating} if
    $\tr(L(\rho)) = 0$ for all $\rho \in \MN{d}$.
    This is the infinitesimal version of trace preservation.
\end{definition}

\begin{definition}[Trace constraint]\label{def:is_trace_constraint}
    \leanok
    \uses{def:generator_decomp}
    A generator decomposition $(\phi, \kappa)$ satisfies the
    \emph{trace constraint} if $\phi^*(\Id) = \kappa + \kappa^\dagger$,
    where $\phi^*$ is the Hilbert--Schmidt adjoint.
    This is the condition in \cite[Eq.~(7.20)]{Wolf2012Quantum}.
\end{definition}

\begin{theorem}[Trace constraint implies trace-annihilating]
    \label{thm:traceAnnihilating_of_traceConstraint}
    \leanok
    \uses{def:generator_decomp, def:is_trace_annihilating,
          def:is_trace_constraint}
    If $(\phi, \kappa)$ satisfies $\phi^*(\Id) = \kappa + \kappa^\dagger$,
    then $L(\rho) = \phi(\rho) - \kappa\rho - \rho\kappa^\dagger$ is
    trace-annihilating.
\end{theorem}

\begin{proof}
    \leanok
    Let $\phi(\rho) = \sum_i K_i \rho K_i^\dagger$.
    By the cyclic property of trace,
    $\tr(L(\rho))
    = \tr((\sum_i K_i^\dagger K_i)\rho)
    - \tr(\kappa\rho)
    - \tr(\kappa^\dagger\rho)
    = \tr((\kappa + \kappa^\dagger)\rho)
    - \tr(\kappa\rho) - \tr(\kappa^\dagger\rho) = 0$.
\end{proof}

\subsection{The Lindblad form}

\begin{definition}[Lindblad form]\label{def:lindblad_form}
    \leanok
    A \emph{Lindblad form} consists of a Hermitian matrix $H = H^\dagger$
    (the Hamiltonian) and a family of matrices $\{L_j\}_{j=0}^{r-1}$
    (the Lindblad operators), defining the linear map
    \[
        L(\rho)
        = i[\rho, H]
        + \sum_{j} \Bigl(        L_j \rho L_j^\dagger
            - \tfrac{1}{2}\{L_j^\dagger L_j, \rho\}_+
          \Bigr),
    \]
    where $[A,B] = AB - BA$ and $\{A,B\}_+ = AB + BA$.
    This is \cite[Eq.~(7.21)]{Wolf2012Quantum}.
\end{definition}

\begin{theorem}[Lindblad form is trace-annihilating]
    \label{thm:lindbladForm_isTraceAnnihilating}
    \leanok
    \uses{def:lindblad_form, def:is_trace_annihilating}
    Every Lindblad form defines a trace-annihilating linear map.
\end{theorem}

\begin{proof}
    \leanok
    The Hamiltonian part satisfies
    $\tr(i[\rho,H]) = 0$ by the cyclic property.
    Each dissipator term satisfies
    $\tr(L_j \rho L_j^\dagger) = \tr(L_j^\dagger L_j \rho)
    = \tr(\rho L_j^\dagger L_j)$,
    so the three terms sum to zero.
\end{proof}

\begin{theorem}[Lindblad form equals generator decomposition]
    \label{thm:lindbladForm_eq_generatorDecomp}
    \leanok
    \uses{def:lindblad_form, def:generator_decomp}
    A Lindblad form with Hamiltonian $H$ and operators $\{L_j\}$ defines
    the same linear map as the generator decomposition $(\phi, \kappa)$
    with $\phi(\rho) = \sum_j L_j \rho L_j^\dagger$ and
    $\kappa = iH + \frac{1}{2}\sum_j L_j^\dagger L_j$.
    This is \cite[Eq.~(7.24)]{Wolf2012Quantum}.
\end{theorem}

\begin{proof}
    \leanok
    Compute $\kappa^\dagger = -iH + \frac{1}{2}\sum_j L_j^\dagger L_j$
    using $H^\dagger = H$ and $(A^\dagger A)^\dagger = A^\dagger A$.
    Setting $S = \sum_j L_j^\dagger L_j$, expand
    $\phi(\rho) - \kappa\rho - \rho\kappa^\dagger
    = \sum_j L_j\rho L_j^\dagger - iH\rho - \tfrac{1}{2}S\rho
      + i\rho H - \tfrac{1}{2}\rho S$,
    which is $i[\rho,H] + \sum_j(L_j\rho L_j^\dagger
    - \tfrac{1}{2}L_j^\dagger L_j\rho - \tfrac{1}{2}\rho L_j^\dagger L_j)$.
\end{proof}

\begin{theorem}[Lindblad form is CCP]
    \label{thm:lindbladForm_isCCP}
    \leanok
    \uses{def:lindblad_form, def:is_ccp,
          thm:lindbladForm_eq_generatorDecomp}
    Every Lindblad form defines a conditionally completely positive map.
\end{theorem}

\begin{proof}
    \leanok
    \uses{thm:lindbladForm_eq_generatorDecomp}
    By Theorem~\ref{thm:lindbladForm_eq_generatorDecomp}, the Lindblad
    form equals a generator decomposition with $\phi$ CP.
\end{proof}

\subsection{Prop.~7.2: Characterization of conditional complete positivity}

\begin{theorem}[Prop.~7.2, direction 2$\to$1]
    \label{thm:choi_projected_implies_ccp}
    \notready
    \uses{def:is_ccp}
    If $L$ is Hermiticity-preserving and
    $P((L \otimes \id)(\ketbra{\Omega}{\Omega}))P \ge 0$
    where $P = \Id - \ketbra{\Omega}{\Omega}$,
    then $L$ is CCP.
    This is \cite[Prop.~7.2]{Wolf2012Quantum}.
\end{theorem}

\begin{proof}
    \notready
    Decompose the Hermitian Choi matrix as
    $\tau_L = Q - \ket{\psi}\bra{\Omega} - \ket{\Omega}\bra{\psi}$
    with $Q \ge 0$ supported on $\ket{\Omega}^\perp$.
    The Choi--Jamiolkowski isomorphism applied to $Q$ gives
    the CP map $\phi$, and
    $(\kappa \otimes \Id)\ket{\Omega} = \ket{\psi}$ defines $\kappa$.
\end{proof}

\subsection{Prop.~7.3: CP semigroups and CCP generators}

\begin{theorem}[Prop.~7.3: CP semigroup implies CCP generator]
    \label{thm:cp_semigroup_implies_ccp}
    \notready
    \uses{def:is_ccp, def:exp_semigroup_ch12}
    If $e^{tL}$ is CP for all $t \ge 0$, then $L$ is CCP.
    This is \cite[Prop.~7.3]{Wolf2012Quantum}.
\end{theorem}

\begin{proof}
    \notready
    From $(e^{tL} \otimes \id)(\ketbra{\Omega}{\Omega}) \ge 0$,
    expand at infinitesimal $t$ and project onto $P$ to get
    $P(L\otimes\id)(\ketbra{\Omega}{\Omega})P \ge 0$.
\end{proof}

\begin{theorem}[Prop.~7.3: CCP generator implies CP semigroup]
    \label{thm:ccp_implies_cp_semigroup}
    \notready
    \uses{def:is_ccp, def:exp_semigroup_ch12}
    If $L$ is CCP, then $e^{tL}$ is CP for all $t \ge 0$.
    This is \cite[Prop.~7.3]{Wolf2012Quantum}.
\end{theorem}

\begin{proof}
    \notready
    Write $L = \phi + \phi_\kappa$ with
    $\phi_\kappa(\rho) = -\kappa\rho - \rho\kappa^\dagger$.
    By Lie--Trotter,
    $e^{tL} = \lim_{n \to \infty}(e^{t\phi/n}e^{t\phi_\kappa/n})^n$.
    Both factors are CP:
    $e^{t\phi/n}$ by Taylor expansion of a CP map, and
    $e^{t\phi_\kappa/n}(\rho) = K\rho K^\dagger$ with $K = e^{-t\kappa/n}$.
\end{proof}

\begin{theorem}[Prop.~7.3: CP semigroup iff CCP generator]
    \label{thm:cp_semigroup_iff_ccp}
    \leanok
    \uses{thm:cp_semigroup_implies_ccp, thm:ccp_implies_cp_semigroup}
    $e^{tL}$ is CP for all $t \ge 0$ if and only if $L$ is CCP.
\end{theorem}

\begin{proof}
    \leanok
    Theorems~\ref{thm:cp_semigroup_implies_ccp}
    and~\ref{thm:ccp_implies_cp_semigroup}.
\end{proof}

\subsection{Prop.~7.4: Freedom in generator representation}

\begin{theorem}[Traceless Kraus operators exist]
    \label{thm:exists_traceless_kraus_shift}
    \leanok
    \uses{def:generator_decomp}
    Given any Kraus operators $\{K_j\}$, there exist $c_j \in \C$
    such that $K_j' = K_j + c_j \Id$ is traceless.
    This is part of \cite[Prop.~7.4]{Wolf2012Quantum}.
\end{theorem}

\begin{proof}
    \leanok
    Set $c_j = -\tr(K_j)/d$.
\end{proof}

\begin{theorem}[Shift invariance of generators]
    \label{thm:generator_shift_invariance}
    \leanok
    \uses{def:generator_decomp}
    If $K_j' = K_j + c_j \Id$ and $\kappa'$ is adjusted per
    \cite[Eq.~(7.19)]{Wolf2012Quantum},
    then $(\phi', \kappa')$ and $(\phi, \kappa)$ define the same generator.
\end{theorem}

\begin{proof}
    \leanok
    Direct algebraic expansion.
\end{proof}

\subsection{Thm.~7.1: The GKSL/Lindblad theorem}

\begin{definition}[GKSL generator]\label{def:is_gksl_generator}
    \leanok
    \uses{def:exp_semigroup_ch12}
    A linear map $L : \MN{d} \to \MN{d}$ is a \emph{GKSL generator}
    if $e^{tL}$ is a quantum channel (CPTP) for all $t \ge 0$.
\end{definition}

\begin{theorem}[Thm.~7.1, Form~(i)]
    \label{thm:gksl_of_generatorDecomp}
    \notready
    \uses{def:is_gksl_generator, def:generator_decomp,
          def:is_trace_constraint, thm:ccp_implies_cp_semigroup}
    If $L(\rho) = \phi(\rho) - \kappa\rho - \rho\kappa^\dagger$ with
    $\phi$ CP and $\phi^*(\Id) = \kappa + \kappa^\dagger$,
    then $L$ is a GKSL generator.
    This is \cite[Thm.~7.1, Eq.~(7.20)]{Wolf2012Quantum}.
\end{theorem}

\begin{proof}
    \notready
    CP follows from CCP via Theorem~\ref{thm:ccp_implies_cp_semigroup}.
    TP follows from
    Theorem~\ref{thm:traceAnnihilating_of_traceConstraint}.
\end{proof}

\begin{theorem}[Thm.~7.1: GKSL iff CCP + trace-annihilating]
    \label{thm:gksl_iff_ccp_ta}
    \notready
    \uses{def:is_gksl_generator, def:is_ccp, def:is_trace_annihilating}
    $L$ is a GKSL generator if and only if $L$ is CCP and
    trace-annihilating.
\end{theorem}

\begin{proof}
    \notready
    Combines Props.~7.2--7.3 with the trace constraint.
\end{proof}

\begin{theorem}[Thm.~7.1: GKSL iff Lindblad form]
    \label{thm:gksl_iff_lindblad}
    \notready
    \uses{def:is_gksl_generator, def:lindblad_form,
          thm:gksl_iff_ccp_ta, thm:lindbladForm_eq_generatorDecomp,
          thm:exists_traceless_kraus_shift}
    $L$ is a GKSL generator if and only if
    \[
        L(\rho)
        = i[\rho, H]
        + \sum_j \Bigl(        L_j \rho L_j^\dagger
            - \tfrac{1}{2}\{L_j^\dagger L_j, \rho\}_+
          \Bigr)
    \]
    with $H = H^\dagger$.
    This is \cite[Thm.~7.1]{Wolf2012Quantum}.
\end{theorem}

\begin{proof}
    \notready
    \uses{thm:gksl_iff_ccp_ta, thm:lindbladForm_eq_generatorDecomp,
          thm:exists_traceless_kraus_shift}
    Forward: by CCP, write $L(\rho) = \phi(\rho) - \kappa\rho
    - \rho\kappa^\dagger$; take traceless Kraus operators
    (Theorem~\ref{thm:exists_traceless_kraus_shift}) and set
    $\kappa = iH + \frac{1}{2}\phi^*(\Id)$
    (Theorem~\ref{thm:lindbladForm_eq_generatorDecomp}).
    Reverse: Theorem~\ref{thm:gksl_iff_ccp_ta}.
\end{proof}

\subsection{Kossakowski matrix form}

\begin{definition}[Kossakowski form]\label{def:kossakowski_form}
    \leanok
    A \emph{Kossakowski form} consists of $H = H^\dagger$,
    a basis $\{F_k\}_{k=0}^{d^2-2}$ of traceless matrices,
    and a PSD matrix $C \in \MN{d^2-1}$, defining
    \[
        L(\rho)
        = i[\rho, H]
        + \tfrac{1}{2}\sum_{k,l} C_{lk}\bigl(        [F_k, \rho F_l^\dagger] + [F_k\rho, F_l^\dagger]
          \bigr).
    \]
    This is \cite[Eq.~(7.23)]{Wolf2012Quantum}.
\end{definition}

\begin{theorem}[Kossakowski form iff Lindblad form]
    \label{thm:kossakowski_iff_lindblad}
    \leanok
    \uses{def:kossakowski_form, def:lindblad_form}
    A linear map admits a Kossakowski form iff it admits a Lindblad form.
\end{theorem}

\begin{proof}
    \leanok
    Diagonalize $C = M^\dagger M$ and set $L_j = \sum_k M_{jk} F_k$.
    Conversely, expand traceless $L_j = \sum_k M_{jk} F_k$ and set
    $C = M^\dagger M \ge 0$.
\end{proof}